% =============================
% Tractable Circuit Zoo - Conceptual Definitions
% Auto-generated from database.json
% =============================
\documentclass[11pt]{article}
\usepackage[margin=1in]{geometry}
\usepackage{amsmath, amssymb, amsthm}
\usepackage{mathtools}
\usepackage{xparse}
\usepackage{hyperref}
\usepackage{natbib}
\hypersetup{colorlinks=true, linkcolor=blue, citecolor=blue, urlcolor=blue}



\NewDocumentCommand{\langref}{m g}{\textbf{#1\IfNoValueF{#2}{#2}}}
\NewDocumentCommand{\langfam}{m m g}{\textbf{#1$_{#2}$\IfNoValueF{#3}{#3}}}
\newcommand{\thislang}{\mathcal{L}}
\newcommand{\poly}{polynomial}
\newcommand{\nopolyunknownquasi}{not polynomial, quasipolynomial unknown}
\newcommand{\nopolyquasi}{not polynomial, quasipolynomial}
\newcommand{\unknownpolyquasi}{polynomial unknown, quasipolynomial}
\newcommand{\unknownboth}{unknown}
\newcommand{\noquasi}{not quasipolynomial}
\newcommand{\CO}{CO}
\newcommand{\VA}{VA}
\newcommand{\CE}{CE}
\newcommand{\IM}{IM}
\newcommand{\EQ}{EQ}
\newcommand{\SE}{SE}
\newcommand{\CT}{CT}
\newcommand{\ME}{ME}
\newcommand{\CD}{CD}
\newcommand{\FO}{FO}
\newcommand{\SFO}{SFO}
\newcommand{\NOTC}{$\neg$C}
\newcommand{\ANDC}{$\wedge$C}
\newcommand{\ANDBC}{$\wedge$BC}
\newcommand{\ORC}{$\vee$C}
\newcommand{\ORBC}{$\vee$BC}
\newcommand{\compilespoly}[2]{#1 compiles polynomially to #2}
\newcommand{\compilesquasi}[2]{#1 compiles quasipolynomially to #2}
\newcommand{\nocompilespoly}[2]{#1 does not compile polynomially to #2}
\newcommand{\nocompilesquasi}[2]{#1 does not compile quasipolynomially to #2}
\newcommand{\supportspoly}[2]{#1 supports #2 in polynomial time}
\newcommand{\supportsquasi}[2]{#1 supports #2 in quasipolynomial time}
\newcommand{\nosupportspoly}[2]{#2 is not supported by #1 in polynomial time}
\newcommand{\nosupportsquasi}[2]{#2 is not supported by #1 in quasipolynomial time}
\newcommand{\isin}[1]{#1}
\newcommand{\shortname}[1]{\textbf{Short name:} #1\par}
\newcommand{\fullname}[1]{\textbf{Full name:} #1\par}
\newcommand{\source}[1]{\textbf{Source:} #1\par}
\newcommand{\target}[1]{\textbf{Target:} #1\par}
\newcommand{\language}[1]{\textbf{Language:} #1\par}
\newcommand{\operation}[1]{\textbf{Operation:} #1\par}
\newcommand{\status}[1]{\textbf{Status:} #1\par}
\newcommand{\selector}[1]{\textbf{Selector:} #1\par}
\newcommand{\assuming}[1]{\textbf{Assuming:} #1\par}
\newcommand{\titlefield}[1]{\textbf{Title:} #1\par}
\let\titlemacro\title
\renewcommand{\title}[1]{\titlefield{#1}}
\newenvironment{description}{\par\noindent\ignorespaces}{\par}
\newenvironment{language}{}{}
\newenvironment{concept}{}{}
\newenvironment{succinctnessclaim}{}{}
\newenvironment{queryclaim}{}{}
\newenvironment{transformationclaim}{}{}
\newenvironment{batchclaim}{}{}
\titlemacro{Tractable Circuit Zoo: Conceptual Definitions}
\date{\today}
\begin{document}
\maketitle

\begin{concept}
\title{Representation Language}
\begin{description}
A \emph{representation language} $L$ is a set of Boolean formulas sharing a common syntactic structure, endowed with the standard Boolean semantics and a size measure $|\cdot|$, usually simply the count of all variables and connectives in the formula. \citep{Darwiche_2002}
\end{description}
\end{concept}

\begin{concept}
\title{Language Family}
\begin{description}
A \emph{language family} is a parameterized collection $\mathcal{F} = \{L_\theta\}_{\theta \in \Theta}$ of representation languages. The parameter $\theta$ may be, for example, a variable order or a vtree.

For example, \langfam{OBDD}{<} is parameterized by a variable order, and \langref{OBDD} is the union over all such parameters: \langref{OBDD} $= \bigcup_<$ \langfam{OBDD}{<}. Distinguishing family-level from language-level claims is central in this project. \citep{Darwiche_2002,Amarilli_2018}

We overload notation in two controlled ways. When the text says ``for each fixed order $<$'' or ``for each fixed vtree $T$,'' notation such as \langfam{OBDD}{<} or \langfam{SDD}{T} denotes the corresponding member language. When no particular parameter is fixed and a family-level claim is intended, the same notation denotes the whole parameterized family, e.g.\ $\{$\langfam{OBDD}{<}$\}_{<}$ or $\{$\langfam{SDD}{T}$\}_{T}$.
\end{description}
\end{concept}

\begin{concept}
\title{Succinctness Between Languages}
\begin{description}
There are two intuitive definitions of succinctness: a \emph{time} requirement, i.e. that an efficient compilation algorithm exists between two languages; or the \emph{space} requirement, that formulas remain ``small'' in the new language. Since, for any computation, runtime is lower-bounded by output size, the latter definition is weaker.

We fix the convention that succinctness between representation languages is a size-only relation, with no requirement that a conversion algorithm exist \citep{Gogic_1995}. However, to the best of our knowledge, all known proofs of succinctness use a compilation algorithm to establish the size relationship; and all proofs that \emph{no} time-efficient compilation algorithm exists use space complexity bounds. An example to the contrary would be highly interesting.

Formally, let $L_1$ and $L_2$ be two languages. $L_1$ is \emph{at least as succinct as} $L_2$, written $L_1 \leq_p L_2$, if there exists a polynomial $p$ such that for every $\alpha \in L_2$ there exists an equivalent $\beta \in L_1$ with $|\beta| \leq p(|\alpha|)$.

We say $L_1$ \emph{polynomially compiles into} $L_2$ if there exists a polynomial-time algorithm that converts any $\alpha \in L_1$ into an equivalent $\beta \in L_2$.

Since all claims in this project work with either definition, we will sometimes use them as if they were interchangeable.

We additionally define quasipolynomial succinctness identically, except we require a quasipolynomial function $q(n) \in n^{log^{O(1)}(n)}$ satisfying the above properties and use the notation $\leq_q$. \citep{Darwiche_2002}
\end{description}
\end{concept}

\begin{concept}
\title{Succinctness Between Families}
\begin{description}
Let $\mathcal{F}_1$ and $\mathcal{F}_2$ be language families. We write $\mathcal{F}_1 \leq_p \mathcal{F}_2$ if for every $L_2 \in \mathcal{F}_2$ there exists $L_1 \in \mathcal{F}_1$ such that $L_1 \leq_p L_2$.

Likewise, we say $\mathcal{F}_1$ polynomially compiles into $\mathcal{F}_2$ if for every $L_1 \in \mathcal{F}_1$ there exists $L_2 \in \mathcal{F}_2$ such that $L_1$ polynomially compiles into $L_2$.

The same convention applies between a language $L$ and a family $\mathcal{F}$ by identifying $L$ with the singleton family $\{L\}$. \citep{Darwiche_2002,Gogic_1995}
\end{description}
\end{concept}

\begin{concept}
\title{Tractability of an Operation}
\begin{description}
An operation $\mathcal{O}$ is \emph{tractable} on a representation language $L$ if there exists a polynomial-time (resp. quasipolynomial-time) algorithm that performs $\mathcal{O}$ on any input formula(s) in $L$. Following \citet{Selman_1996}, the point is to compile a propositional theory off-line into a language where target operations are tractable. \citep{Darwiche_2002}
\end{description}
\end{concept}

\begin{concept}
\title{Query}
\begin{description}
A \emph{query} takes a formula $\alpha \in L$ (and possibly auxiliary inputs) and returns a Boolean answer or a numerical value without modifying $\alpha$.

The standard query operations in this project are Consistency, Validity, Clausal Entailment, Implicant, Equivalence, Model Counting, Model Enumeration, and Sentential Entailment. \citep{Darwiche_2002}
\end{description}
\end{concept}

\begin{concept}
\title{Transformation}
\begin{description}
A \emph{transformation} maps one or more formulas from $\alpha_1,\alpha_2,..., \alpha_m\ \in L$, where $m$ may be arbitrarily large, to a new formula $\alpha_{m+1} \in L$.

The standard transformations in this project are Conditioning, Forgetting, Singleton Forgetting, Conjunction, Bounded Conjunction, Disjunction, Bounded Disjunction, and Negation. \citep{Darwiche_2002}
\end{description}
\end{concept}

\begin{concept}
\title{Consistency (CO)}
\begin{description}
Given a representation $\alpha$, decide whether $\alpha$ is satisfiable, i.e., whether there exists an assignment $X$ such that $\alpha(X)=\top$. \citep{Darwiche_2002}
\end{description}
\end{concept}

\begin{concept}
\title{Validity (VA)}
\begin{description}
Given a representation $\alpha$, decide whether $\alpha$ is valid, i.e., $\alpha \equiv \top$. \citep{Darwiche_2002}
\end{description}
\end{concept}

\begin{concept}
\title{Clausal Entailment (CE)}
\begin{description}
Given a representation $\alpha$ and a clause $c$, decide whether $\alpha \models c$ (equivalently, whether $\alpha \wedge \neg c$ is unsatisfiable). \citep{Darwiche_2002}
\end{description}
\end{concept}

\begin{concept}
\title{Implicant (IM)}
\begin{description}
Given a representation $\alpha$ and a term $t$, decide whether $t \models \alpha$ (equivalently, whether $t \wedge \neg \alpha$ is unsatisfiable). \citep{Darwiche_2002}
\end{description}
\end{concept}

\begin{concept}
\title{Equivalence (EQ)}
\begin{description}
Given two representations $\alpha$ and $\beta$, decide whether $\alpha \equiv \beta$. \citep{Darwiche_2002}
\end{description}
\end{concept}

\begin{concept}
\title{Model Counting (CT)}
\begin{description}
Given a representation $\alpha$, compute the number of satisfying assignments: $|\{X:\alpha(X)=\top\}|$. \citep{Darwiche_2002}
\end{description}
\end{concept}

\begin{concept}
\title{Model Enumeration (ME)}
\begin{description}
Given a representation $\alpha$, enumerate all satisfying assignments, i.e., output $\{X:\alpha(X)=\top\}$, in time polynomial in the input size and in the total size of the output. \citep{Darwiche_2002}
\end{description}
\end{concept}

\begin{concept}
\title{Sentential Entailment (SE)}
\begin{description}
Given representations $\alpha$ and $\beta$, decide whether $\alpha \models \beta$. \citep{Darwiche_2002}
\end{description}
\end{concept}

\begin{concept}
\title{Conditioning (CD)}
\begin{description}
Given a representation $\alpha$ and a consistent term $\gamma$, compute the conditioned representation $\alpha|_{\gamma}$ obtained by substituting the assignments in $\gamma$. \citep{Darwiche_2002}
\end{description}
\end{concept}

\begin{concept}
\title{Forgetting (FO)}
\begin{description}
Given a representation $\alpha$ and a set of variables $Y$, compute $\exists Y.\alpha$ (existentially quantifying the variables in $Y$). \citep{Darwiche_2002}
\end{description}
\end{concept}

\begin{concept}
\title{Singleton Forgetting (SFO)}
\begin{description}
Given a representation $\alpha$ and a variable $x$, compute $\exists x.\alpha = (\alpha|_{x=0}) \vee (\alpha|_{x=1})$. \citep{Darwiche_2002}
\end{description}
\end{concept}

\begin{concept}
\title{Conjunction ($\wedge$C)}
\begin{description}
Given representations $\alpha_1,\ldots,\alpha_m$, compute $\bigwedge_{i=1}^m \alpha_i$. \citep{Darwiche_2002}
\end{description}
\end{concept}

\begin{concept}
\title{Bounded Conjunction ($\wedge$BC)}
\begin{description}
Given two representations $\alpha_1,\alpha_2$, compute $\alpha_1 \wedge \alpha_2$. \citep{Darwiche_2002}
\end{description}
\end{concept}

\begin{concept}
\title{Disjunction ($\vee$C)}
\begin{description}
Given representations $\alpha_1,\ldots,\alpha_m$, compute $\bigvee_{i=1}^m \alpha_i$. \citep{Darwiche_2002}
\end{description}
\end{concept}

\begin{concept}
\title{Bounded Disjunction ($\vee$BC)}
\begin{description}
Given two representations $\alpha_1,\alpha_2$, compute $\alpha_1 \vee \alpha_2$. \citep{Darwiche_2002}
\end{description}
\end{concept}

\begin{concept}
\title{Negation ($\neg$C)}
\begin{description}
Given a representation $\alpha$, compute $\neg \alpha$. \citep{Darwiche_2002}
\end{description}
\end{concept}
\bibliographystyle{plainnat}
\bibliography{refs}
\end{document}
